\begin{proof}
    Since all generators in $\Sigma_{\CatHQC}$ are endomorphic, then \cref{Thm:DefCtrl} is applicable to $\CatHQC$.
    Let $\{ \tau_n: \CatHQC( n, n ) \to \CatHQC( n + 1, n + 1 ) \}_{n \in \mathbb{N}}$ be the family of functions defined by the equation $\tau_n( U ) = \scalebox{0.65}{\circimg{tikz-cache/dagger__term1.pdf}} \odot U$.
    The six conditions hold as follows.
    \begin{enumerate}
    \item If $n \in \mathbb{N}$, $\tau_n( 1_n ) \approx_E 1 \boxtimes 1_n = 1_{n+1}$ by \HqcCtrlId.
    \item If $U: n \to n$ and $V: n \to n$, then $\tau_n( V \circ U ) = \tau_n( V ) \circ \tau_n( U )$ by \HqcCtrlComp.
    \item If $U: n \to n$ and $( U, V ) \in E$, then $\tau_n( U ) \approx_E \tau_n( V )$ by \HQCCtrlCong.
    \item Let $U: n \to n$.
          Then the following derivation holds. 
          \begin{align*}
              \scalebox{0.65}{\circimg{tikz-cache/proofs__ctrl_tensorator_1.pdf}}
              \overunderset{\HQCCtrlCong}{\HqcCtrlPair}{\approx}
              \scalebox{0.65}{\circimg{tikz-cache/proofs__ctrl_tensorator_2.pdf}}
              \overset{\HqcCtrlComp}{\approx}
              \scalebox{0.65}{\circimg{tikz-cache/proofs__ctrl_tensorator_3.pdf}}
              \overset{\HqcCtrlPair}{\approx}
              \scalebox{0.65}{\circimg{tikz-cache/proofs__ctrl_tensorator_4.pdf}}
          \end{align*}
          Then $\tau_{n+1}( U \boxtimes 1 ) \approx_E \tau_n( U ) \boxtimes 1$.
    \item Let $U: n \to n$. Then the following implication holds.
          Since $\sigma^\dagger = \sigma$, then the following derivation also holds. 
          \begin{align*}
            \scalebox{0.63}{\circimg{tikz-cache/proofs__ctrl_sym_9.pdf}}
            \overset{\HqcCtrlSwap}{\approx}&
            \scalebox{0.63}{\circimg{tikz-cache/proofs__ctrl_sym_10.pdf}}
            \overunderset{\HQCCtrlCong}{\HqcCtrlSwap}{\approx}
            \scalebox{0.63}{\circimg{tikz-cache/proofs__ctrl_sym_5.pdf}}
            \overset{\HqcCtrlDag}{\approx}
            \scalebox{0.63}{\circimg{tikz-cache/proofs__ctrl_sym_6.pdf}}
            \\
            =\;\;\;&
            \scalebox{0.63}{\circimg{tikz-cache/proofs__ctrl_sym_7.pdf}}
            \overset{\HqcCtrlSwap}{\approx}
            \scalebox{0.63}{\circimg{tikz-cache/proofs__ctrl_sym_8.pdf}}
          \end{align*}
          Then $\tau_{n+1}( \tau_n( U ) ) \approx_E ( \sigma \boxtimes 1_n ) \circ \tau_{n+1}( \tau_n( U ) ) \circ ( \sigma \boxtimes 1_n )$.
    \item Let $U: n + 2 \to n + 2$ and $0 \le k \le n$.
          Since $1_k^\dagger = 1_k$, $1_{n-k}^\dagger = 1_{n-k}$, and $\sigma^\dagger = \sigma$, then it follows that $\gamma_{n,k}^\dagger = ( 1_k \boxtimes \sigma \boxtimes 1_{n-k} )^\dagger = 1_k^\dagger \boxtimes \sigma^\dagger \boxtimes 1_{n-k}^\dagger = \gamma_{n,k}$.
          Since $\gamma_{n,k}^\dagger = \gamma_{n,k}$, then it follows by \HqcCtrlDag that $\tau_{n+2}( \gamma_{n,k} \circ U \circ \gamma_{n,k} ) \approx_E (1 \boxtimes \gamma_{n,k} ) \circ \tau_{n+2}( U ) \circ ( 1 \boxtimes \gamma_{n,k} )$.
    \end{enumerate}
    Since conditions (1) and (2) hold, then by \cref{Thm:DefCtrl} there exists a unique ordinary functor $F: \CatHQC \to \CatHQC / E$ such that $F_0( n ) = n + 1$ and $F( U: n \to n ) = \pi_E( \tau_n( U ) )$.
    Since condition (3) also holds, then by \cref{Thm:DefCtrl} exists a unique ordinary functor $C: \CatHQC / E \to \CatHQC / E$ such that $H \circ \pi_E = F$.
    Then $C( \pi_E( U ) ) = F( U ) = \pi_E( \scalebox{0.65}{\circimg{tikz-cache/dagger__term1.pdf}} \odot U )$ for each $U \in \CatHQC( n, n )$.
    Since conditions (4) through to (6) hold, then $C$ is a control functor.
    Since every morphism in $\CatHQC / E$ is unitary with respect to $( \overline{\dagger} )$ by \cref{Thm:HQCInv}, then $C$ is a dagger control functor with respect to $( \overline{\dagger} )$ by \cref{Lem:UnitaryCtrl}.
    
    It remains to be shown that $C$ is a conjugated control functor.
    Let $X \in ( \CatHQC / E )( n, n )$ and $Y \in ( \CatHQC / E )( n, n )$.
    Then there exists $U \in \CatHQC( n, n )$ and $V \in \CatHQC( n, n )$ such that $\pi_E( U ) = X$ and $\pi_E( V ) = Y$.
    Then $( 1 \boxtimes U^\dagger ) \circ ( \scalebox{0.65}{\circimg{tikz-cache/dagger__term1.pdf}} \odot V ) \circ ( 1 \boxtimes U ) \approx_E ( \scalebox{0.65}{\circimg{tikz-cache/dagger__term1.pdf}} \odot ( U^\dagger \circ V \circ U ) )$ by $\HqcCtrlDag$.
    Then the following equation holds.
    \begin{align*}
        \left( 1 \boxtimes X^{\overline{\dagger}} \right) \circ C( \pi_E( V ) ) \circ ( 1 \boxtimes X )
        &=
        \left( 1 \boxtimes \pi_E( U )^{\overline{\dagger}} \right) \circ C( \pi_E( V ) ) \circ ( 1 \boxtimes \pi_E( U ) )
        \\
        &=
        \left( 1 \boxtimes \pi_E\left( U^{\dagger} \right) \right) \circ C( \pi_E( V ) ) \circ ( 1 \boxtimes \pi_E( U ) )
        &\text{(by \cref{Thm:HQCInv})}
        \\
        &=
        \pi_E\left( \scalebox{0.65}{\circimg{tikz-cache/proofs__ctrl_conj_eqn_term2.pdf}} \right) \circ C( \pi_E( V ) ) \circ \pi_E\left( \scalebox{0.65}{\circimg{tikz-cache/proofs__ctrl_conj_eqn_term1.pdf}} \right)
        \\
        &=
        \pi_E\left( \scalebox{0.65}{\circimg{tikz-cache/proofs__ctrl_conj_eqn_term2.pdf}} \right) \circ \pi_E\left( \scalebox{0.65}{\circimg{tikz-cache/proofs__ctrl_conj_eqn_term3.pdf}} \right) \circ \pi_E\left( \scalebox{0.65}{\circimg{tikz-cache/proofs__ctrl_conj_eqn_term1.pdf}} \right)
        &\text{(by \cref{Thm:HQCCtrl})}
        \\
        &=
        \pi_E\left( \scalebox{0.65}{\circimg{tikz-cache/proofs__ctrl_conj_eqn_term4.pdf}} \right)
        \\
        &=
        \pi_E\left( \scalebox{0.65}{\circimg{tikz-cache/proofs__ctrl_conj_eqn_term5.pdf}} \right)
        &\text{(by $\approx_E$)}
        \\
        &=
        C \left( \pi_E\left( \scalebox{0.65}{\circimg{tikz-cache/proofs__ctrl_conj_eqn_term6.pdf}} \right) \right)
        &\text{(by \cref{Thm:HQCCtrl})}
        \\
        &=
        C\left( \pi_E\left( U^\dagger \right) \circ \pi_E( V ) \circ \pi_E( U ) \right)
        \\
        &=
        C\left( \pi_E( U )^{\overline{\dagger}} \circ \pi_E( V ) \circ \pi_E( U ) \right)
        \\
        &=
        C\left( X^{\overline{\dagger}} \circ Y \circ X \right)
        &\text{(by \cref{Thm:HQCInv})}
    \end{align*}
    Since $X$ and $Y$ were arbitrary, then $C( - )$ is conjugated.
\end{proof}
